\section{Zwei Zufallsvariablen}

Wenn zwei Zufallsvariablen $X$ und $Y$ gegeben sind, ist die Verbundverteilungsfunktion durch

$$
F_{XY}(x,y) = P(X \leq x, Y \leq y)
$$

definiert. Es gilt

$$
F_{XY}(-\infty, y) = 0, \quad F_{XY}(x, -\infty) = 0 \quad \text{und} \quad F_{XY}(\infty, \infty) = 1.
$$

Die Verbundverteilungsdichtefunktion ist dann durch

$$
p_{XY}(x,y) = \frac{\partial^2 F_{XY}(x,y)}{\partial x \, \partial y}
$$

gegeben. Somit gilt auch

$$
F_{XY}(x,y) = \int_{-\infty}^{x} \int_{-\infty}^{y} p_{XY}(\xi, \mu) \, \mathrm{d}\xi \, \mathrm{d}\mu.
$$

\subsection{Statistische Unabhängigkeit}

Für zwei statistisch unabhängige Zufallsvariablen $X$ und $Y$ gilt:

$$
F_{XY}(x,y) = F_X(x) \cdot F_Y(y)
$$

und

$$
p_{XY}(x,y) = p_X(x) \cdot p_Y(y).
$$

\subsection{Randverteilungen}

Die univariate Verteilungsdichtefunktion der Zufallsvariablen $X$ und $Y$ lässt sich mit

$$
p_X(x) = \int_{-\infty}^{\infty} p_{XY}(x,y)\,\mathrm{d}y \quad \text{und} \quad p_Y(y) = \int_{-\infty}^{\infty} p_{XY}(x,y)\,\mathrm{d}x
$$

bestimmen. Diese werden Randverteilungen genannt. Für die Verteilungsfunktion ergibt sich

$$
F_X(x) = F_{XY}(x, \infty) \quad \text{und} \quad F_Y(y) = F_{XY}(\infty, y).
$$

\begin{beispiel}[Bivariate Normalverteilung]
Zwei Zufallsvariablen $X$ und $Y$ sind gemeinsam normalverteilt, wenn ihre Verbundverteilungsdichtefunktion durch

$$
p_{XY}(x,y) = \frac{1}{2\pi\sigma_X\sigma_Y\sqrt{1-r_{XY}^2}} \exp\!\left(-\frac{1}{2(1-r_{XY}^2)}\left[\frac{(x-\mu_X)^2}{\sigma_X^2} - 2r_{XY}\frac{(x-\mu_X)(y-\mu_Y)}{\sigma_X\sigma_Y} + \frac{(y-\mu_Y)^2}{\sigma_Y^2}\right]\right)
$$

gegeben ist. Hierbei sind $\mu_X$ und $\mu_Y$ die Mittelwerte, $\sigma_X^2$ und $\sigma_Y^2$ die Varianz und $r_{XY}$ der Korrelationskoeffizient

$$
r_{XY} = \frac{\mathrm{E}\{(X-\mu_X)(Y-\mu_Y)\}}{\sigma_X\sigma_Y}
$$

der Zufallsvariablen $X$ und $Y$. Es gilt $|r_{XY}| \leq 1$.

% Grafik: 3D-Oberfläche und Konturplot von p_XY(x,y) für μ_X=0, μ_Y=0, σ_X²=1, σ_Y²=1, r_XY=0

% Grafik: 3D-Oberfläche und Konturplot von p_XY(x,y) für μ_X=1, μ_Y=1, σ_X²=1, σ_Y²=1, r_XY=0

% Grafik: 3D-Oberfläche und Konturplot von p_XY(x,y) für μ_X=0, μ_Y=0, σ_X²=1, σ_Y²=0.5, r_XY=0

% Grafik: 3D-Oberfläche und Konturplot von p_XY(x,y) für μ_X=0, μ_Y=0, σ_X²=1, σ_Y²=1, r_XY=0.5
\end{beispiel}

\subsection{Momente der bivariaten Verteilungsdichte}

Die Momente zweier Zufallsvariablen $X$ und $Y$ mit der VDF $p_{XY}(x,y)$ lassen sich mit dem folgenden Theorem ermitteln:

$$
\mathrm{E}\{g(X,Y)\} = \int_{-\infty}^{\infty} \int_{-\infty}^{\infty} g(x,y)\,p_{XY}(x,y)\,\mathrm{d}x\,\mathrm{d}y.
$$

Hierbei ist $g(X,Y)$ eine beliebige Funktion der Zufallsvariablen $X$ und $Y$.

Speziell gilt im Fall $Z = g(X,Y) = aX + bY$

$$
\mathrm{E}\{aX + bY\} = a\,\mathrm{E}\{X\} + b\,\mathrm{E}\{Y\}.
$$

Die Erwartungswertbildung ist eine lineare Operation.

\subsection{Linearität des Erwartungswertes}

Für zwei Zufallsvariablen $X$ und $Y$ gilt:

$$
\mathrm{E}\{aX + bY\} = a\,\mathrm{E}\{X\} + b\,\mathrm{E}\{Y\} \quad \text{für alle } a, b \in \mathbb{R}
$$

\subsection{Kovarianz, Korrelation und Orthogonalität}

Die Kovarianz von zwei Zufallsvariablen $X$ und $Y$ ist zu

$$
C = \mathrm{E}\{(X - \mu_X)(Y - \mu_Y)\}
$$

definiert. Es gilt $C = \mathrm{E}\{XY\} - \mu_X\mu_Y$.

Der Korrelationskoeffizient ist zu

$$
r_{XY} = \frac{\mathrm{E}\{(X - \mu_X)(Y - \mu_Y)\}}{\sigma_X\sigma_Y}
$$

definiert. Es gilt $|r_{XY}| \leq 1$.

Zwei Zufallsvariablen werden orthogonal genannt, wenn

$$
\mathrm{E}\{XY\} = \int_{-\infty}^{\infty} \int_{-\infty}^{\infty} x\,y\,p_{XY}(x,y)\,\mathrm{d}x\,\mathrm{d}y = 0
$$

gilt.

\subsection{Eine Funktion von zwei Zufallsvariablen}

Es sei eine Funktion von zwei Zufallsvariablen $X$ und $Y$ gegeben:

$$
Z = g(X, Y)
$$

Die Verteilungsfunktion $F_Z(z)$ ist dann durch

$$
F_Z(z) = P(Z \leq z) = P(g(X,Y) \leq z) = \iint_{D_z} p_{XY}(x,y)\,\mathrm{d}x\,\mathrm{d}y
$$

gegeben, wobei $D_z$ das nicht unbedingt zusammenhängende Gebiet in der $xy$-Ebene bezeichnet, für das $g(x, y) \leq z$ gilt.

\begin{beispiel}
Für $Z = X + Y$ gilt:

% Grafik: Gebiet z ≤ x+y in der xy-Ebene (Halbebene unterhalb der Linie y = z-x)

$$
F_Z(z) = \int_{-\infty}^{\infty} \int_{-\infty}^{z-x} p_{XY}(x,y)\,\mathrm{d}y\,\mathrm{d}x
$$

$$
p_Z(z) = \int_{-\infty}^{\infty} p_{XY}(x, z-x)\,\mathrm{d}x
$$

und im Fall unabhängiger Zufallsvariablen $X$ und $Y$

$$
p_Z(z) = \int_{-\infty}^{\infty} p_X(x)\,p_Y(z-x)\,\mathrm{d}x
$$
\end{beispiel}

\subsection{Zwei Funktionen von zwei Zufallsvariablen}

Gegeben sind zwei Zufallsvariablen $X$ und $Y$ und zwei Funktionen

$$
Z = g(X, Y) \quad \text{und} \quad W = h(X, Y).
$$

Die Verbundverteilungsdichtefunktion der Zufallsvariablen $Z$ und $W$ ist dann durch

$$
p_{ZW}(z,w) = \frac{p_{XY}(x_1, y_1)}{|\det(J(x_1, y_1))|} + \ldots + \frac{p_{XY}(x_n, y_n)}{|\det(J(x_n, y_n))|}
$$

gegeben, wobei die $x_i$ die Lösungen des Gleichungssystems $g(x_i, y_i) = z$, $h(x_i, y_i) = w$ sind und $J$ die Jakobi-Matrix

$$
J(x,y) = \begin{pmatrix} \dfrac{\partial g}{\partial x} & \dfrac{\partial g}{\partial y} \\[6pt] \dfrac{\partial h}{\partial x} & \dfrac{\partial h}{\partial y} \end{pmatrix}
$$

bezeichnet. Hierbei ist $\det(J(x_i, y_i))$ die Determinante dieser Matrix.

\subsection{Statistische Mittelwerte und Zeitmittelwerte}

% Grafik: Drei diskrete Realisierungen X_1[k], X_2[k], X_3[k] eines Zufallsprozesses als Impulsfolgen

\subsection{Ausblick: Zufallsprozesse}

Zufällige diskrete Signale lassen sich nun als Folge von Zufallsvariablen $X[k]$ definieren. Dabei kann für jede Zufallsvariable eine Verteilungsdichtefunktion und entsprechende statistische Mittelwerte und Momente angegeben werden. Zudem wird die statistische Abhängigkeit der Zufallsvariablen mit Verbundverteilungsdichtefunktionen beschrieben. Für die Modellbildung sind besonders

\begin{itemize}
  \item stationäre Zufallsprozesse (alle $X[k]$ weisen die selbe Verteilungsdichtefunktion auf) und
  \item ergodische Zufallsprozesse (die statistischen Mittelwerte entsprechen den durch Mittelung über eine Folge $X[k]$ gewonnenen empirischen Mittelwerten).
\end{itemize}


